\section{条件熵、联合熵与链式法则}


熵描述单个随机变量的不确定性；条件熵描述“已知 $Y$ 之后，$X$ 还剩多少不确定性”。本节的核心公式是链式法则

\[
H(X,Y)=H(Y)+H(X\mid Y),
\]

以及“条件不会增加熵”：

\[
H(X\mid Y)\le H(X).
\]

\textbf{本节用到的记号：} $P_{X,Y}(x,y)=P(X=x,Y=y)$；$P_{X\mid Y}(x\mid y)=P(X=x\mid Y=y)$；$H(X\mid Y=y)$ 为给定 $Y=y$ 的条件熵；$H(X\mid Y)$ 为对这些条件熵按 $Y$ 的分布取平均。

\subsection{条件熵的定义}


给定 $Y=y$，条件分布为 $P_{X\mid Y=y}$，其熵为

\[
H(X\mid Y=y)
=\sum_xP_{X\mid Y}(x\mid y)\log\frac1{P_{X\mid Y}(x\mid y)}.
\]

再对 $Y$ 的取值按 $P_Y(y)$ 加权平均，得到\textbf{条件熵}：

\[
H(X\mid Y)
=\sum_yP_Y(y)H(X\mid Y=y).
\]

\textbf{注意.} $H(X\mid Y=y)$ 是一个数；$H(X\mid Y)$ 是一个数（不是随机变量）——它已经把 $Y$ 平均掉了。若把 $H(X\mid Y)$ 写成随机变量，应先定义 $H(X\mid Y)(\omega)$，再取期望；本课程不需要这样写。

\textbf{特例.} 若 $Y$ 是常数，则 $H(X\mid Y)=H(X)$；若 $Y=X$，则 $H(X\mid Y)=0$。

\subsection{联合熵}


两个随机变量的联合熵为

\[
H(X,Y)
=\sum_{x,y}P_{X,Y}(x,y)\log\frac1{P_{X,Y}(x,y)}.
\]

它衡量的是“同时观察 $X,Y$ 时总的平均不确定性”。

\subsection{链式法则}


\textbf{定理.}

\[
\boxed{
H(X,Y)=H(Y)+H(X\mid Y)
}
\]

对称地也有 $H(X,Y)=H(X)+H(Y\mid X)$。

\textbf{证明.} 把联合分布写成 $P_{X,Y}(x,y)=P_Y(y)P_{X\mid Y}(x\mid y)$：

\[
\begin{aligned}
H(X,Y)
&=\sum_{x,y}P_{X,Y}(x,y)\log\frac1{P_{X,Y}(x,y)}\\
&=\sum_{x,y}P_{X,Y}(x,y)\left[
\log\frac1{P_Y(y)}+\log\frac1{P_{X\mid Y}(x\mid y)}
\right]\\
&=\sum_yP_Y(y)\log\frac1{P_Y(y)}
+\sum_{x,y}P_{X,Y}(x,y)\log\frac1{P_{X\mid Y}(x\mid y)}\\
&=H(Y)+H(X\mid Y).
\end{aligned}
\]

最后一项正是条件熵的定义。

\textbf{多变量链式法则.} 对 $X_1,\dots,X_n$，

\[
H(X_1,\dots,X_n)
=H(X_1)+H(X_2\mid X_1)+\cdots
+H(X_n\mid X_1,\dots,X_{n-1}).
\]

证明：反复用两变量链式法则，把最后一个变量单独提出来：

\[
H(X_1,\dots,X_n)=H(X_1,\dots,X_{n-1})+H(X_n\mid X_1,\dots,X_{n-1}),
\]

再对前 $n-1$ 个变量归纳。

\subsection{条件熵的非负性与“条件不增加熵”}


\textbf{非负性.}

\[
H(X\mid Y)\ge0.
\]

因为它是若干条件分布的熵的加权平均，而每个熵都非负。

\textbf{条件不增加熵.}

\[
\boxed{
H(X\mid Y)\le H(X)
}
\]

等号成立当且仅当 $X,Y$ 独立。

\textbf{证明（用 KL/对数和）.} 计算差值：

\[
\begin{aligned}
H(X)-H(X\mid Y)
&=\sum_xP_X(x)\log\frac1{P_X(x)}
-\sum_{x,y}P_{X,Y}(x,y)\log\frac1{P_{X\mid Y}(x\mid y)}\\
&=\sum_{x,y}P_{X,Y}(x,y)\log\frac{P_{X\mid Y}(x\mid y)}{P_X(x)}\\
&=\sum_{x,y}P_{X,Y}(x,y)\log\frac{P_{X,Y}(x,y)}{P_X(x)P_Y(y)}.
\end{aligned}
\]

这是 $P_{X,Y}$ 与“乘积分布”$P_XP_Y$ 之间的 KL 散度（见第 4 节）。由 KL 非负性可得差值 $\ge0$；等号当且仅当 $P_{X,Y}=P_XP_Y$，即 $X,Y$ 独立。

（如果还没学 KL，也可以直接用 Jensen：$\log$ 凹，$\sum p\log y\le\log\sum py$。）

\textbf{联合熵的上界.} 由链式法则和条件熵不等式，

\[
H(X,Y)=H(X)+H(Y\mid X)\le H(X)+H(Y),
\]

等号当且仅当 $X,Y$ 独立。多变量情形同样有

\[
H(X_1,\dots,X_n)\le\sum_iH(X_i).
\]

\subsection{条件熵为零的含义}


\[
H(X\mid Y)=0
\]

当且仅当：给定 $Y$ 后，$X$ 没有不确定性，即 $X$ 是 $Y$ 的函数（存在可测函数 $f$ 使 $X=f(Y)$ 几乎必然）。

这与“$H(X,Y)=H(Y)$”等价：联合的不确定性全在 $Y$ 上。

\subsection{熵的凹性（补证）}


上一节说 $H$ 是分布的凹函数。这里用条件熵给一个干净的证明。

设 $P=\lambda P_0+(1-\lambda)P_1$，$\lambda\in[0,1]$。构造随机变量 $Z\sim\mathrm{Bern}(\lambda)$，并令

\[
X\mid Z=0\sim P_0,\qquad X\mid Z=1\sim P_1.
\]

则 $X$ 的边缘分布正好是 $P$。由链式法则，

\[
H(X,Z)=H(X)+H(Z\mid X)=H(Z)+H(X\mid Z).
\]

其中：

\begin{itemize}
\item $H(Z)=h(\lambda)$；
\item $H(X\mid Z)=\lambda H(P_0)+(1-\lambda)H(P_1)$；
\item 由“条件不增加熵”，$H(Z\mid X)\le H(Z)$。
\end{itemize}


于是

\[
\begin{aligned}
H(P)=H(X)
&=H(Z)+H(X\mid Z)-H(Z\mid X)\\
&\ge H(X\mid Z)\\
&=\lambda H(P_0)+(1-\lambda)H(P_1).
\end{aligned}
\]

这正是凹性。等号当且仅当 $H(Z\mid X)=H(Z)$；在混合模型中，这通常要求 $P_0=P_1$（否则知道 $X$ 会改变 $Z$ 的后验）。

\textbf{结论.}

\[
H(\lambda P_0+(1-\lambda)P_1)
\ge\lambda H(P_0)+(1-\lambda)H(P_1).
\]

\subsection{常见错误}


\begin{itemize}
\item 把 $H(X\mid Y)$ 与 $H(X\mid Y=y)$ 混为一谈：前者是平均后的数。
\item 把链式法则写成 $H(X,Y)=H(X)+H(Y)$；只有独立时才能这样拆。
\item 认为 $H(X\mid Y)\le H(X)$ 的等号总是成立；等号要求独立。
\item 忘记“条件熵”是“平均条件熵”，不是“某个 $y$ 下的条件熵”。
\end{itemize}
